\documentclass[sigconf,anonymous,review]{acmart}

%% ---- placeins provides \FloatBarrier if float packing needs a nudge.
\usepackage{placeins}

%% ---- author-year citations; the body emits \citep / \citet from
%%      pandoc-style markers in the canonical source.
\citestyle{acmauthoryear}

%% ---- neutral bibliographic stubs (manuscript under review) -----------
\acmConference[Preprint]{Manuscript under review}{2026}{}
\acmYear{2026}
\settopmatter{printacmref=false}
\setcopyright{none}
\renewcommand\footnotetextcopyrightpermission[1]{}
\pagestyle{plain}

%% ---- ACM CCS concepts stub (real values land at camera-ready) --------
\begin{CCSXML}
<ccs2012>
  <concept>
    <concept_id>10002978.10003014.10003017</concept_id>
    <concept_desc>Security and privacy~Intrusion/anomaly detection and malware mitigation</concept_desc>
    <concept_significance>500</concept_significance>
  </concept>
</ccs2012>
\end{CCSXML}
\ccsdesc[500]{Security and privacy~Intrusion/anomaly detection and malware mitigation}

\keywords{deterministic verification, read-depth robustness, adversarial evasion, pre-registration, multi-agent LLM security, LLM-judge audit}

\title{Reading Deeper Re-Acquires the Attack: A Pre-Registered Read-Depth Robustness Trilemma for Deterministic Threat Verifiers}

%% ---- acmart anonymous mode requires at least one author block --------
\author{Anonymous Author(s)}
\affiliation{%
  \institution{Submission under double-blind review}
  \country{}
}
\email{anonymous@paper4.example}

\begin{document}

\begin{abstract}
Deterministic verification is the standard defense against framing attacks on multi-agent LLM security pipelines: replacing an LLM adjudicator with fixed rules removes the interpretive surface a meaning-preserving paraphrase can exploit. The defense buys stability at a cost that is rarely measured --- a deterministic rule reads less of the evidence than an LLM would, so detection power and framing-robustness pull against each other. This paper asks whether a deterministic verifier can be given \textit{deeper reading} to recover detection power without re-acquiring the framing vulnerability it was deployed to escape. We define a five-rung read-depth ladder --- from field-presence checks up through canonicalization to full LLM reading --- and, on a pre-registered adversarial corpus (Adaptive Corpus C, 8 scenarios), measure every rung against three properties at once: detection (Youden's \textit{J}), framing-robustness (a noise-gated majority-flip-rate under meaning-preserving operators), and data-integrity (not inheriting attacker-controllable structure). With the decision bands frozen \textit{before} the run, the good corner of the frontier is empty: on the deployment-realistic cumulative view every rung caps at \textit{J} = 0.25, and the only rung that occupies the corner in the standalone view --- deterministic canonicalization, at \textit{J} = 0.75 and zero framing-flip --- is shown to do so by construction. A blind out-of-vocabulary adversary defeats that lone escape on 2 of 4 malign string scenarios; the two scenarios that turn on a named indicator (lsass, procdump) resist evasion in both arms, and white-box access does not expand the evadable set. Two independent judge families (GPT-4o, Gemini), under calibration controls, corroborate that the evasions preserve threat meaning rather than reflecting a lenient judge. The result is an empirical, pre-registered \textit{read-depth robustness trilemma}: removing the LLM from the reading seat does not remove the attack --- it relocates it, and climbing back toward deep reading re-opens the door. All findings are reproducible from closed, pre-registered artifacts.

\end{abstract}
\maketitle


\section{Introduction}
Multi-agent LLM pipelines now make security and judgment decisions in production, and a recurring design move hardens them against adversarial input: replace the LLM that \textit{adjudicates} the evidence with a deterministic, non-LLM verifier. The motivation is well established. LLM-integrated applications are vulnerable to indirect prompt injection, where adversarial content rides in on the data the model reads rather than the operator's prompt \citep{greshake-2023}, and the threat broadens in multi-agent settings where any participating agent's evidence stream can be perturbed \citep{guo-2024}. The ARES program established the defensive half of this picture: a four-rule deterministic Skeptic, composed in front of a deterministic adjudicator, matches a full-LLM Skeptic's accuracy on a framing-injection corpus at zero LLM-call exposure \citep{gmys-casiano-2026b}, and the resulting verdicts are byte-stable even where the same pipeline's LLM-sourced explanation surface drifts \citep{gmys-casiano-2026c}. Removing the LLM from the reading seat works because a fixed rule has no interpretive surface for a paraphrase to capture.

That defense, however, buys its stability with shallowness. A deterministic rule reads less of the evidence than an LLM would: it can check that a field is present, or match a literal token, but it cannot weigh the meaning of an analyst's narrative the way a language model can. The framing-robust verifier and the deep-reading verifier sit at opposite ends of a spectrum, and the natural engineering question is whether the gap can be closed --- whether a deterministic verifier can be given \textit{deeper reading}, recovering the detection power of an LLM, without dragging back the framing vulnerability it was deployed to escape.

This paper answers that question by measurement, and the answer is a tension we argue is fundamental rather than incidental. We define a five-rung \textit{read-depth ladder}, from a field-presence check (\texttt{v1\_field}) through structured, lexical, and canonicalizing rules (\texttt{v2\_structured}, \texttt{v2\_lexical}, \texttt{v2\_canonical}) up to full LLM reading (\texttt{llm\_semantic}), and we evaluate every rung against three properties simultaneously: \textbf{detection} (does it catch real threats?), \textbf{framing-robustness} (does a meaning-preserving rewording flip its verdict?), and \textbf{data-integrity} (does it stay honest, or does it "detect" only by trusting surface structure the attacker himself controls?). Our central claim is a \textit{read-depth robustness trilemma}: across this ladder, on a pre-registered adversarial corpus, no rung is simultaneously deep-reading, framing-robust, and data-integrity-clean. The deepest reader (the LLM) detects the most but is talked out of a correct verdict by a single reworded sentence; the shallow deterministic rungs resist rewording but either miss too much or "catch" threats only by matching attacker-controllable structure. The trilemma is sharpest in the gap between two views of the frontier: read independently, a rung can look like it occupies the good corner, but stacked into a deployable pipeline its detection collapses to the cap set by the most error-prone rung beneath it.

We make three contributions. \textbf{First}, a pre-registered read-depth robustness frontier: with the framing-robust and high-detection bands frozen before the measurement run, the good corner is empty --- on the cumulative (deployment-realistic) view, all five rungs cap at Youden's \textit{J} = 0.25. \textbf{Second}, an adversarial refutation of the lone apparent escape: deterministic canonicalization sits in the good corner standalone (\textit{J} = 0.75 at zero framing-flip), but we flagged in advance that its robustness might be an artifact of undoing only the disguises it was built for, and a blind out-of-vocabulary LLM adversary --- proposing brand-new disguises the rule had never seen --- defeats it on 2 of 4 malign string scenarios; tellingly, the two scenarios anchored on a named tool (lsass, procdump) resist, because a named indicator cannot be disguised away while still meaning the same thing, and white-box access does not widen the breach. \textbf{Third}, an independent-judge audit and a reusable method: because the evasion verdict rests on an LLM judge calling the disguises meaning-preserving, we re-judge them with two other model families (GPT-4o and Gemini) under calibration controls, which corroborate the evasions as genuine threats; the recipe --- pre-registered bands, an LLM-proposes / code-disposes adversary, and an independent audit with calibration controls --- transfers to any LLM-judge-reliant evaluation.

We are explicit about the honesty boundary of the claim. This is an empirical, pre-registered frontier result on one corpus, one ladder, and one pipeline model family --- not an impossibility proof; "trilemma" is the organizing claim, not a theorem (Section 10). Pre-registration is what converts the empty corner from a curve fit into a credible result: the bands and decision rule were committed to the audit record before the deciding run, so the goalposts could not move after the data arrived. The remainder of the paper proceeds as follows. Section 2 positions the work against the framing-injection, deterministic-verification, and detector-evasion literatures. Section 3 compresses the ARES verifier and defines the read-depth ladder. Section 4 specifies the pre-registered frontier methodology. Sections 5, 6, and 7 present the three findings. Section 8 generalizes the method, Section 9 discusses the relocation thesis, and Sections 10 and 11 state limitations and future work.


\section{Background}
The architectural pattern under study composes one or more LLM agents with a non-LLM adjudication layer that consumes their outputs and emits a verdict. \citet{guo-2024} survey the rapid growth of this multi-agent family between 2023 and 2024, across debate, role specialization, and pipeline orchestration; such pipelines now appear in security triage, fraud analysis, and compliance, wherever a deterministic step is layered over LLM-generated hypotheses to provide reproducibility. The adversarial pressure these systems face is indirect prompt injection: \citet{greshake-2023} established that LLM-integrated applications can be compromised by adversarial content embedded in the data an agent reads rather than the prompt the operator issues, and in a multi-agent pipeline that content can be perturbed within the evidence stream of any participating agent. The defensive reframe the ARES program advances is that deterministic verification over a closed, schema-checked evidence set does not eliminate such semantic attacks but converts them into data-integrity attacks against the evidence itself \citep{gmys-casiano-2026b}, and that the verdict's decision surface can be made byte-stable even while its LLM-sourced explanation surface drifts \citep{gmys-casiano-2026c}. The closed-world commitment that makes this possible --- agents reason only over a bound, immutable fact set --- is the classical posture formalized by \citet{reiter-1978}.

The reframe has an unmeasured cost. A deterministic verifier earns its framing-robustness precisely by \textit{not reading} the interpretive surface a paraphrase would capture; the more of the evidence's meaning it declines to read, the less it can detect. This is the gap the present paper measures: between a shallow rule that cannot be talked out of a verdict because it barely reads, and a deep reader that detects subtle threats but re-exposes the interpretive surface. Framing the question as a spectrum of \textit{read depth} --- rather than a binary deterministic-versus-LLM choice --- is what lets us locate where, along that spectrum, robustness and detection become jointly unattainable.

Two literatures supply the measurement vocabulary. The first is robustness-frontier and pre-registration practice in machine-learning evaluation: characterizing a method by where it sits in a tradeoff space rather than a single headline number \citep{tsipras-2019}, and committing decision rules before the deciding run so that a favorable result cannot be manufactured by post-hoc band selection \citep{albanie-2022}. Pre-registration is the mechanism that converts our central "empty corner" claim from a curve fit into a credible result, and we adopt it as a first-class methodological commitment rather than a reporting afterthought. The second is the evasion of pattern-matching detectors: a rule with a fixed vocabulary has a natural adversarial bypass through synonym substitution and out-of-vocabulary rephrasing that preserves meaning while defeating the literal match \citep{jin-2020}. This is exactly the failure mode our adversarial probe (Section 6) is designed to elicit against the one deterministic rung that otherwise looks robust.

Finally, the data-integrity property we measure connects to the faithfulness literature in explainable AI. \citet{jacovi-goldberg-2020} formalize faithfulness as the requirement that an explanation represent the underlying computation rather than merely appear plausible. The data-integrity leg of our trilemma is a stability analogue of that requirement: a verifier has \textit{low} data-integrity when its detection rests on surface structure the attacker controls --- when, in effect, it "explains" its catch by pointing at a pattern the adversary can rewrite at will. A rule that detects only by matching attacker-controllable structure is brittle in exactly the way an unfaithful explanation is misleading, and Section 4 operationalizes this leg through the inheritance of a component's false-positive rate when readers are stacked.


\section{The ARES Verifier and the Read-Depth Ladder}
\begin{figure*}[!htbp]
\centering
\includegraphics[width=\textwidth]{../figures/fig_1.pdf}
\caption{The read-depth ladder: five rungs from v1\_field through v2\_structured, v2\_lexical, and v2\_canonical up to llm\_semantic. Each rung reads more of the evidence value than the one below it, raising detection potential but enlarging the attacker-controlled surface.}
\Description{The read-depth ladder: five rungs from v1\_field through v2\_structured, v2\_lexical, and v2\_canonical up to llm\_semantic. Each rung reads more of the evidence value than the one below it, raising detection potential but enlarging the attacker-controlled surface.}
\label{fig:fig_1}
\end{figure*}

The ARES (Adversarial Reasoning Engine System) pipeline composes three agents in fixed order: an Architect that reads a bound \texttt{EvidencePacket} and emits a threat hypothesis with a cited fact set; a Skeptic that produces an analytical response; and a deterministic OracleJudge that consumes both through a closed-form decision table and emits a \texttt{Verdict}, never issuing an LLM call. Two commitments make the pipeline auditable and are carried unchanged from the preceding ARES papers \citep{gmys-casiano-2026b, gmys-casiano-2026c}: the \texttt{EvidencePacket} is a frozen, SHA256-verified set of typed \texttt{Fact} records that no agent may extend, and the pipeline reasons closed-world over that bound set. The object of study in this paper is the Skeptic position --- specifically, how much of each \texttt{Fact}'s content the Skeptic reads when deciding whether the evidence is malign.

The deterministic Light Skeptic is a pure-Python rule engine. We generalize it into a \textbf{read-depth ladder} of five rungs that read progressively more of the evidence value content. \texttt{v1\_field} reads only field \textit{presence} --- whether a high-threat field is populated at all --- and nothing of its value. \texttt{v2\_structured} adds structural logic over those fields: a high-threat field combined with a kill-chain stage and no authorization marker. \texttt{v2\_lexical} reads the field \textit{values} through fixed regular expressions --- an executable in a user-writable path, credential-access tooling, an ineffective-patch marker. \texttt{v2\_canonical} first canonicalizes a value (normalizing synonyms and surface variants) and then matches, so that \texttt{update.exe} in \texttt{\%TEMP\%} and \texttt{update.exe} in \texttt{C:\textbackslash{}Temp\textbackslash{}} collapse to the same signal. At the top, \texttt{llm\_semantic} reads the evidence by meaning through a full LLM call, the deepest reader on the ladder. The four deterministic rungs and the LLM anchor are registered as \texttt{DETERMINISTIC\_TIERS} and \texttt{LADDER\_ORDER} in the codebase, which fix the rung definitions against silent drift.

The ladder makes the central tension concrete. Each rung up reads more of the value content the rung below it ignored, and reading more of the value is what raises detection --- but the value is exactly the surface an attacker controls. A field-presence rule cannot be fooled by \textit{what} the value says because it never reads the value; a regex rule can be fooled by a value the attacker phrases outside the pattern; a canonicalizing rule can be fooled by a synonym its canonicalizer does not know; and an LLM reader, which reads everything, can be steered by a meaning-preserving paraphrase of the whole. Climbing the ladder trades one weakness for another. The remainder of the paper measures that trade across all five rungs at once.


\section{Methodology: A Pre-Registered Frontier}

\subsection{Adaptive Corpus C}
The measurement runs on \textbf{Adaptive Corpus C}, eight stratified scenarios built to probe read depth specifically. Four are \textit{value-borne malign} scenarios whose threat lives in the field value, not the field structure, across three strata: two lexical (\texttt{RDF-M-LEX-001}, anchored on the credential-dumping target \texttt{lsass}; \texttt{RDF-M-LEX-002}, an executable dropped in a temp path), one synonym (\texttt{RDF-M-SYN-001}, a loader described in rephrasable prose), and one patch (\texttt{RDF-M-PATCH-001}, anchored on \texttt{procdump}). The remaining four are negative and positive controls: benign \textit{structural twins} that share a malign scenario's field structure but carry benign values (so a structure-only rung that fires on them is revealed as trusting structure rather than detecting threat), a carry-forward false-positive probe built from benign \texttt{.com}/\texttt{.js} indicators, and an \texttt{inject\_authorization} positive control that should flip the structural tier by construction. The malign-versus-benign split is what the detection metric scores; the controls are what expose a rung that "detects" by inheriting structure. The corpus is content-addressed by a digest (\texttt{9401b7188ba790a5}) so every downstream measurement pins the exact scenario set.


\subsection{Three properties}
Each rung is scored on three properties simultaneously. \textbf{P1, detection}, is Youden's \textit{J} = TPR $-$ FPR over the malign scenarios versus their benign twins: a rung scores high only by catching malign evidence \textit{and} sparing benign evidence. \textbf{P2, framing-robustness}, is \textit{X}\_sem, a noise-gated majority-flip rate: we apply meaning-preserving framing operators (a prefix and a suffix paraphrase) to a malign scenario, resample the rung's verdict under each, and count a flip only when the perturbed majority verdict differs from the baseline majority \textit{and} a permutation test rejects sampling noise at \textit{p} < 0.05. A framing-robust rung has \textit{X}\_sem near zero. \textbf{P3, data-integrity}, asks whether a rung's detection survives being stacked with the rungs beneath it rather than resting on structure a shallower rung already mis-fires on; Section 4.3 makes this precise as an inheritance property rather than a third free axis.


\subsection{Standalone versus cumulative views}
We read the frontier in two views. The \textbf{standalone} view measures each rung in isolation --- its own TPR, FPR, and \textit{J} as if it were the only reader. The \textbf{cumulative} view measures each rung \textit{stacked on all the rungs beneath it}, as it would actually deploy: a production verifier does not discard its cheaper checks when it adds a deeper one. The two views diverge because false positives accumulate. If a shallow rung fires on a benign structural twin, every deeper rung stacked above it inherits that false positive, and the stack's \textit{J} is capped by the worst component FPR below it. This inheritance is the operational form of the data-integrity leg (P3): a rung whose standalone detection looks strong but whose stacked detection collapses is one whose apparent strength depended on the very structure a shallower rung already trusts. P3 therefore enters the frontier not as a third independent axis but as the penalty the cumulative view imposes --- a deliberate two-axis operationalization of a three-property claim, stated openly here and revisited as a limitation in Section 10.


\subsection{Frozen bands and the decision rule}
Before the deciding run, we committed two bands and a decision rule to the pre-registration record. A rung is \textbf{framing-robust} iff \textit{X}\_sem $\leq$ 0.10, and \textbf{high-detection} iff cumulative \textit{J} $\geq$ 0.50. The \textbf{good corner} of the (\textit{X}\_sem, \textit{J}) frontier is the region satisfying both. The pre-registered decision rule: the trilemma is \textit{supported} if the good corner is empty --- no rung is simultaneously framing-robust and high-detection on the cumulative view --- and \textit{falsified} if any rung occupies it. Freezing the bands and the rule before the run is the methodological core of the paper: it removes the post-hoc degree of freedom by which an "empty corner" could otherwise be manufactured, and it commits us in advance to a result that could have gone the other way. The bands are mirrored between the pre-registration document and the code constants that compute the verdict, and an always-on test asserts the two never diverge.


\subsection{Live anchor}
The deterministic rungs are computed offline, but the deciding run anchors the ladder against a live LLM tier: the \texttt{llm\_semantic} rung and the framing-flip probe are measured with a real model (Sonnet 4, K = 20 resamples) so that the LLM rung's detection and framing-susceptibility are empirical rather than assumed. The malign verdict for a rung is its threat-confirmed reading of a scenario; the framing-flip metric resamples that reading under the meaning-preserving operators. Section 5 reports the resulting frontier; Sections 6 and 7 then stress the one rung that survives the frozen bands.


\section{Finding 1: The Good Corner Is Empty}
\begin{figure*}[!htbp]
\centering
\includegraphics[width=\textwidth]{../figures/fig_2.pdf}
\caption{The read-depth frontier (the central result): framing-flip rate X\_sem versus Youden's J, standalone (left) versus cumulative (right). The good corner is occupied standalone by deterministic canonicalization but empty on the deployment-realistic cumulative view, where every rung caps at J = 0.25.}
\Description{The read-depth frontier (the central result): framing-flip rate X\_sem versus Youden's J, standalone (left) versus cumulative (right). The good corner is occupied standalone by deterministic canonicalization but empty on the deployment-realistic cumulative view, where every rung caps at J = 0.25.}
\label{fig:fig_2}
\end{figure*}


\subsection{Result statement}
On the deciding tier-4 run (Sonnet 4, K = 20 resamples, Adaptive Corpus C digest \texttt{9401b7188ba790a5}, \$3.23), the good corner of the (\textit{X}\_sem, cumulative \textit{J}) frontier is empty: no rung of the read-depth ladder is simultaneously framing-robust (\textit{X}\_sem $\leq$ 0.10) and high-detection (cumulative \textit{J} $\geq$ 0.50). Under the pre-registered decision rule (\S{}4.4), this \textbf{supports} the trilemma. The result is not that deep reading fails to detect, nor that shallow reading fails to resist framing; it is that the two properties are never jointly present in a deployable configuration. The frontier divides sharply by view, and the division is the finding.


\subsection{The standalone ladder: detection rises with depth}
Read in isolation, the ladder behaves exactly as the read-depth intuition predicts --- detection climbs monotonically as each rung reads more of the evidence value. Standalone Youden's \textit{J} rises \texttt{0.00 $\rightarrow$ 0.25 $\rightarrow$ 0.50 $\rightarrow$ 0.75} across \texttt{v1\_field}, \texttt{v2\_structured}, \texttt{v2\_lexical}, \texttt{v2\_canonical}, and the LLM rung also reaches \textit{J} = 0.75. The field-presence rung detects nothing it can distinguish from benign (TPR 0.0); the structured rung catches every malign scenario (TPR 1.0) but at a punishing standalone FPR of 0.75; the lexical rung trades some recall for precision (TPR 0.75, FPR 0.25); and canonicalization recovers full recall at that lower false-positive rate (TPR 1.0, FPR 0.25, \textit{J} = 0.75). Taken at face value, the standalone view says deeper reading simply wins.


\subsection{The cumulative collapse: the data-integrity leg}
The deployment-realistic cumulative view erases that story. Stacked as a production verifier actually runs them --- each rung added on top of, not in place of, the rungs beneath it --- every rung from \texttt{v2\_structured} upward caps at \textit{J} = 0.25, and no rung exceeds it. The cause is false-positive inheritance: the structured tier fires on three of four benign controls (standalone FPR 0.75), and because a stack confirms a threat if \textit{any} component does, that 0.75 false-positive rate propagates into every deeper rung's cumulative score. Canonicalization's clean standalone \textit{J} = 0.75 becomes a cumulative \textit{J} = 0.25 the moment it sits above a structural tier that cries wolf. This gap between the standalone and cumulative views is the empirical signature of the data-integrity leg (\S{}4.3): the deeper rungs' apparent detection power is, in deployment, hostage to the structural tier's willingness to trust attacker-shaped structure. The standalone-to-cumulative collapse is the trilemma's teeth.


\subsection{The LLM rung detects but is framing-susceptible}
The deepest reader confirms the other edge of the tension. Standalone, the LLM rung detects well (TPR 0.75, FPR 0.00, \textit{J} = 0.75) --- but it is not framing-robust: its framing-flip rate is \textit{X}\_sem = 0.125 (95\% interval [0.00, 0.375]), placing it \textit{outside} the framing-robust band even before deployment. The single flip is diagnostic. On the synonym scenario \texttt{RDF-M-SYN-001}, a meaning-preserving suffix paraphrase moves the LLM's majority verdict from malign (baseline rate 1.00) to benign (perturbed rate 0.45), and a permutation test rejects sampling noise at \textit{p} = 0.0005; the prefix paraphrase of the same scenario shifts the rate to 0.55 (\textit{p} = 0.0015) without flipping the majority. One reworded sentence is enough to talk the deepest reader out of a correct verdict. The LLM rung thus fails the frontier from the framing axis (\textit{X}\_sem > 0.10), exactly where the deterministic rungs fail from the detection axis (cumulative \textit{J} < 0.50).


\subsection{The lone escape}
One rung resists this account, and only in the standalone view: deterministic canonicalization (\texttt{v2\_canonical}) sits at \textit{X}\_sem = 0.00 and standalone \textit{J} = 0.75 --- inside the good corner. Its framing-robustness is real by construction: a canonicalizer that normalizes before matching cannot be moved by the surface variants it normalizes away, so its framing-flip rate is zero against the operators in the corpus. We flagged this rung in the pre-registration as the named potential non-falsifier: its standalone good-corner position is exactly what a by-construction robustness would produce, and an artifact of undoing only the disguises the canonicalizer already knows is indistinguishable, on this corpus, from genuine robustness. Whether \texttt{v2\_canonical} is a true escape from the trilemma or merely a rule blind to disguises it has not seen is the question Section 6 settles, by turning a blind adversary loose to invent disguises it has never seen.


\section{Finding 2: The Lone Escape Is Adversarially Defeated}
\begin{figure}[!htbp]
\centering
\includegraphics[width=\columnwidth]{../figures/fig_3.pdf}
\caption{Out-of-vocabulary adversarial evasion per scenario: canonical flip counts (of K = 8) for each malign scenario, black-box versus white-box arms. RDF-M-LEX-002 and RDF-M-SYN-001 are evaded in both arms; the named-indicator scenarios RDF-M-LEX-001 (lsass) and RDF-M-PATCH-001 (procdump) resist.}
\Description{Out-of-vocabulary adversarial evasion per scenario: canonical flip counts (of K = 8) for each malign scenario, black-box versus white-box arms. RDF-M-LEX-002 and RDF-M-SYN-001 are evaded in both arms; the named-indicator scenarios RDF-M-LEX-001 (lsass) and RDF-M-PATCH-001 (procdump) resist.}
\label{fig:fig_3}
\end{figure}

\begin{figure*}[!htbp]
\centering
\includegraphics[width=\textwidth]{../figures/fig_6.pdf}
\caption{A worked out-of-vocabulary disguise (RDF-M-SYN-001, black-box arm, confirmed by both independents): the original values the canonicalizer matches versus the disguise it misses, with threat meaning preserved.}
\Description{A worked out-of-vocabulary disguise (RDF-M-SYN-001, black-box arm, confirmed by both independents): the original values the canonicalizer matches versus the disguise it misses, with threat meaning preserved.}
\label{fig:fig_6}
\end{figure*}


\subsection{The probe: LLM proposes, code disposes}
To test whether \texttt{v2\_canonical}'s robustness is genuine or by-construction, we attack it with out-of-vocabulary (OOV) disguises --- rewordings that fall outside the canonicalizer's known synonym set. The generator follows an \textit{LLM-proposes, code-disposes} protocol that separates adversarial creativity from adversarial validity. An LLM adversary proposes disguised values for a malign scenario; a deterministic gate then admits a candidate only if it passes two checks: \textbf{skeleton-invariance}, that the disguise leaves the scenario's structural fields and fact identifiers unchanged (so the attack is on the value the canonicalizer reads, not on the evidence schema), and \textbf{novelty}, that the disguise introduces a token outside the canonicalizer's synonym vocabulary (so it is a genuine OOV attack rather than a known variant the rule already folds away). Only validated, novel disguises reach the rung. This design prevents the adversary from "winning" by corrupting the evidence skeleton or by replaying in-vocabulary variants, and it makes the surviving evasions interpretable: each is a meaning-preserving rephrase the canonicalizer genuinely cannot normalize.


\subsection{Two arms, and the result}
We run the probe in two arms --- \textbf{black-box}, where the adversary sees only the scenario, and \textbf{white-box}, where it also sees the canonicalizer's source --- at K = 8 candidates per scenario per arm (32 candidates per arm across the four malign scenarios). The verdict is \textbf{SUPPORTED\_STRONG}: a blind black-box adversary defeats \texttt{v2\_canonical} on 2 of the 4 malign string scenarios (adversarial flip rate per scenario = 0.50 in each arm), evading on \texttt{RDF-M-LEX-002} (the executable-in-temp scenario) and \texttt{RDF-M-SYN-001} (the synonym loader). The per-candidate flip rate is 0.281 black / 0.290 white --- roughly three in ten validated disguises slip past the canonicalizer. The OOV corpus is content-addressed (\texttt{a4ea1d0645152ffa}) and the run cost \$0.106. The headline is not the flip rate but its existence: a rung that scored \textit{X}\_sem = 0.00 against the in-corpus framing operators is moved off a correct verdict by disguises drawn from outside its vocabulary. Its standalone good-corner position was by-construction robustness after all.


\subsection{Named indicators resist}
The two scenarios that resist are as informative as the two that fall. \texttt{RDF-M-LEX-001} (anchored on \texttt{lsass}) and \texttt{RDF-M-PATCH-001} (anchored on \texttt{procdump}) show zero canonical flips in either arm: no validated disguise evades them. The reason is structural to the threat, not to the rule. These scenarios turn on a \textit{named indicator} --- a specific tool whose name \textit{is} the threat signal --- and a named indicator cannot be disguised away while still meaning the same thing. Rephrasing "lsass" into something that is no longer "lsass" does not preserve the threat; it removes it. Where detection rests on a named entity, meaning-preserving evasion has no room to operate; where it rests on rephrasable structure or generic synonyms (\texttt{v2\_lexical}/\texttt{v2\_canonical} on the LEX-002 and SYN-001 scenarios), it does. This split --- named-IOC matching resists, structural and synonym matching fall --- is the most actionable observation in the paper, and Section 9 develops it into a design implication.


\subsection{White-box $\approx$ black-box}
Source access does not widen the breach. The white-box arm, which can read the canonicalizer's normalization tables, evades the same two scenarios as the blind black-box arm and at a near-identical per-candidate rate (0.290 versus 0.281). The evadable boundary therefore tracks the \textit{attacker's ability to construct a meaning-preserving synonym}, not the attacker's knowledge of the rule's vocabulary. This matters for the threat model: the defense gains little from keeping its normalization tables secret, because the adversary does not need them --- natural-language paraphrase capacity, which any capable LLM has, is sufficient. A canonicalizer's robustness cannot be restored by obscurity.


\subsection{First observation and the persisted run}
We note for reproducibility that the evasion was first observed in an earlier run (OOV corpus \texttt{9900b91f707e2ef8}, \$0.1104, per-candidate flip 0.344 black / 0.312 white) whose disguise texts were not persisted and whose meaning-preservation rested on a single judge. The numbers reported above are from the persisted, audited re-run (corpus \texttt{a4ea1d0645152ffa}), which is the canonical artifact and the one every figure and gate in this paper is locked to. The first run is the observation; the second is the record. The shift from "observed once" to "persisted and independently judged" is the subject of Section 7, because the evasion result rests on a claim that the disguises preserve threat meaning --- and that claim was, in the first run, made by a single LLM judge that accepted every disguise it saw.


\section{Finding 3: Independent-Judge Audit}
\begin{figure}[!htbp]
\centering
\includegraphics[width=\columnwidth]{../figures/fig_4.pdf}
\caption{Independent-judge audit: malign verdicts over the 18 evading disguises by judge family (Sonnet, Gemini, GPT-4o), with GPT-4o the stricter judge. Fifteen disguises are confirmed by both independents and three split; all four calibration controls pass, yielding the ROBUST verdict.}
\Description{Independent-judge audit: malign verdicts over the 18 evading disguises by judge family (Sonnet, Gemini, GPT-4o), with GPT-4o the stricter judge. Fifteen disguises are confirmed by both independents and three split; all four calibration controls pass, yielding the ROBUST verdict.}
\label{fig:fig_4}
\end{figure}


\subsection{The leniency threat}
The SUPPORTED\_STRONG verdict has a load-bearing dependency: it holds only if the disguises that evade \texttt{v2\_canonical} genuinely \textit{preserve threat meaning}. If a disguise turns a malign value into something a security analyst would read as benign, then the canonicalizer was right to pass it and the "evasion" is no evasion at all. In the verdict run, meaning-preservation was judged by a single LLM (Sonnet), and that judge accepted all 64 candidates it was shown with zero rejections. A judge that never says no is indistinguishable, from the outside, from a judge that is simply lenient. The evasion result therefore cannot stand on it alone, and we audit it.


\subsection{An independent panel with calibration controls}
The audit re-judges the evading disguises with two model families from different providers --- GPT-4o and Gemini --- that played no part in generating or first-judging them. Critically, we do not trust the independents blind either: each is first run against four \textbf{calibration controls}. Two are positive (the original, undisguised malign values of the evaded scenarios, which a competent judge \textit{must} read as malign) and two are negative (a benign structural twin and a clean baseline, which a competent judge \textit{must} read as benign). Only a panel that passes its controls --- correctly calling the obvious cases --- is trusted to rule on the disputed disguises. This is the direct leniency test: a judge that rubber-stamps everything fails the negative controls.


\subsection{Result: ROBUST}
All four calibration controls pass: both independents correctly call the positive controls malign and the negative controls benign, so neither is a rubber stamp. On the 18 evading disguises, the two independents both confirm 15 as meaning-preserving threats; the remaining 3 split (GPT-4o reads them benign while Gemini and Sonnet read them malign). Under the pre-registered "both-independents" confirmation rule --- a disguise is confirmed only if \textit{both} GPT-4o and Gemini agree it is malign --- the audit verdict is \textbf{ROBUST}: the evasions are corroborated as genuine threats, not artifacts of a lenient judge. The three splits are not noise but signal: they cluster at the meaning-dilution boundary of the synonym scenario \texttt{RDF-M-SYN-001}, where successive rephrasing eventually thins the threat enough that a stricter reader hesitates. GPT-4o is the stricter judge of the three, and the splits mark where "meaning-preserving" shades into "meaning-diluting." The audit cost \$0.0093, bringing total live spend across the evasion and audit to \$0.115.


\subsection{What the audit settles}
The audit converts the evasion from a one-judge observation into a multi-family result. The canonicalizer's robustness was by-construction (Section 6), and the disguises that defeat it are threats that two independent model families, having proven they can tell malign from benign on controls, agree remain malign (Section 7). The lone escape from the trilemma is closed, and closed under scrutiny designed to catch the most likely way the closing could be wrong.


\section{The Method as a Contribution}
\begin{figure*}[!htbp]
\centering
\includegraphics[width=\textwidth]{../figures/fig_5.pdf}
\caption{The method as a transferable contribution: the LLM-proposes / code-disposes pipeline (generate, then a deterministic skeleton-and-novelty gate, then an LLM judge, then verdict), plus the independent-judge audit panel that separates adversarial creativity from adversarial validity.}
\Description{The method as a transferable contribution: the LLM-proposes / code-disposes pipeline (generate, then a deterministic skeleton-and-novelty gate, then an LLM judge, then verdict), plus the independent-judge audit panel that separates adversarial creativity from adversarial validity.}
\label{fig:fig_5}
\end{figure*}

The instrument that produced this result is, we argue, reusable beyond it, and we state it as a secondary contribution: a three-layer recipe for evaluating a robustness claim that an LLM is in a position to flatter. The first layer is \textbf{pre-registration of bands and decision rule}. An "empty corner" or any tradeoff-space result is only as credible as the commitment that fixed its thresholds before the data arrived; freezing \textit{X}\_sem $\leq$ 0.10, cumulative \textit{J} $\geq$ 0.50, and the supported/falsified rule in advance is what lets us claim the corner is empty rather than chosen to be empty. The second layer is the \textbf{LLM-proposes / code-disposes adversary}. Generative breadth and validity are separated: the LLM supplies disguises no fixed enumerator would think of, and a deterministic gate (skeleton-invariance plus OOV-novelty) admits only the ones that are genuine, in-bounds attacks. The gate is what stops the adversary from "winning" by breaking the evidence schema or replaying known variants, and it makes every surviving evasion interpretable. The third layer is the \textbf{independent-judge audit with calibration controls}. When a result depends on an LLM judging meaning, a single judge that never refuses is worthless as evidence; re-judging with independent model families that must first pass positive and negative controls is a reusable leniency-falsification test --- it converts "an LLM said so" into "judges that demonstrably distinguish malign from benign said so." None of the three layers is novel in isolation; their composition is the method, and it transfers to any setting where a deterministic component's robustness, or any LLM-judged evaluation, could be an artifact of the evaluator rather than a property of the system.


\section{Discussion}

\subsection{The relocation thesis}
The through-line of the ARES program is that hardening a multi-agent LLM pipeline against semantic attack does not destroy the attack --- it relocates it, and the present paper measures one full circuit of that relocation. Paper 1 located the failure inside the black box: multi-turn LLM debate degrades calibration rather than improving it \citep{gmys-casiano-2026a}. Paper 2 moved the defense to a deterministic layer and showed that doing so converts semantic attacks into data-integrity attacks on the evidence schema \citep{gmys-casiano-2026b}. Paper 3 showed the relocation again at finer grain: a deterministic adjudicator can hold the decision surface byte-stable while the LLM-sourced explanation surface drifts \citep{gmys-casiano-2026c}. The read-depth trilemma is the next turn of the same screw. Pull the LLM out of the reading seat and the framing attack that flips its verdict (Section 5.4) disappears --- but the detection you lose forces you to read more of the value, and reading the value re-acquires \textit{either} the framing vulnerability (climb back to the LLM rung) \textit{or} a dependence on attacker-controllable structure (the shallow rungs' cumulative collapse, Section 5.3) \textit{or}, at the one rung that seemed to escape both, an out-of-vocabulary evasion (Section 6). The attack is conserved across the defense; only its address changes. This is why we frame the result as a trilemma rather than a defect to be patched: each of the three properties can be bought, but not all three at once, because buying one re-sells another to the attacker.


\subsection{Not all deterministic rules are equally fragile}
The sharpest practical observation is the named-indicator split (Section 6.3). The rungs that fell to OOV evasion detect by matching rephrasable structure or generic synonyms; the rungs that resisted detect by matching a named entity whose name is the threat. This is not a quirk of the corpus --- it is a property of the threat surface. A meaning-preserving rephrase has purchase exactly where the signal is rephrasable, and none where the signal is a proper noun the adversary cannot rename without removing the threat. Deterministic robustness, in other words, is not a single quantity; it is high for indicator-anchored rules and low for structure-anchored ones, and a verifier's place in the trilemma depends on which kind of rule carries its detection.


\subsection{Implications for production verifier design}
Two design consequences follow. First, the cumulative collapse (Section 5.3) is a warning against stacking a high-FPR structural tier beneath deeper readers: the stack inherits the worst false-positive rate below it, so a cheap, trigger-happy early rung silently caps the detection of everything above it. A production ladder should gate its shallow rungs on precision, not just recall, or the depth it pays for never materializes. Second, the named-indicator result suggests where a defensible corner of the trilemma might actually lie: a hybrid rung that anchors on indicator specificity first and invokes semantic reading only as a secondary, gated step might inherit the OOV-resistance of named-IOC matching while recovering some of the detection of deep reading. We did not build that rung, and we do not claim it escapes the trilemma; we name it as the most promising corner-occupant candidate the data points to (Section 11).


\subsection{The honesty dividend of pre-registration}
Finally, a methodological note. The strength of the empty-corner claim comes less from the numbers than from the fact that the numbers were committed to a rule that could have falsified them. A single rung occupying the good corner on the cumulative view would have read as FALSIFIED under the frozen decision rule, and \texttt{v2\_canonical} came close enough --- occupying the corner standalone --- that the result hinged on a view distinction fixed in advance, not chosen after. Presenting a result that survived a pre-registered chance to fail is worth more than a larger effect presented post hoc, and we regard the pre-registration discipline as part of the contribution rather than scaffolding around it.


\section{Limitations}
The claim is bounded, and we state the bounds plainly. (1) \textbf{One corpus, few malign scenarios.} Adaptive Corpus C has eight scenarios and four malign string scenarios; the frontier is measured on a deliberately small, hand-stratified set, and broader scenario distributions are untested. (2) \textbf{One pipeline model family.} The live frontier and framing-flip probe ran on Sonnet 4; cross-model generalization of the frontier itself is untested, and a different base model could place the LLM rung differently on the framing axis. (3) \textbf{One OOV adversary family.} The disguises were proposed by a single adversary model; a second or third adversary family would strengthen the evasion claim and might widen the evadable set. (4) \textbf{A control gap on the synonym scenario.} \texttt{RDF-M-SYN-001} has no same-skeleton benign twin in the corpus, so its leniency control leans on the generic clean baseline, making the audit slightly weaker for the synonym class than for the lexical class. (5) \textbf{A two-axis operationalization of a three-property claim.} The data-integrity leg (P3) is realized as the cumulative-view false-positive penalty rather than as an independent third axis (\S{}4.3); the "trilemma" is thus measured on a two-axis frontier where the third property enters as a cap. (6) \textbf{Empirical, not a theorem.} "Trilemma" is a pre-registered empirical result on this setup, not a proven impossibility; a different ladder, corpus, or adversary could in principle find a corner occupant, and Section 11's hybrid rung is exactly such a candidate.


\section{Future Work}
Four directions follow directly from the limitations. First, \textbf{a benign synonym-class twin} for \texttt{RDF-M-SYN-001}, closing the audit's control gap (\S{}10.4) and letting the synonym class be scored on the same footing as the lexical class. Second, \textbf{additional OOV adversary families} --- proposing disguises with two or three independent models --- to test whether the evadable set is adversary-dependent and to harden the evasion claim against single-adversary artifacts. Third, \textbf{a larger and broader Corpus C}, with more malign/benign pairs across additional threat categories, to test whether the empty corner survives outside the hand-stratified set. Fourth, and most consequential, \textbf{a named-IOC-plus-semantic hybrid rung}: a tier that anchors on indicator specificity first --- where OOV evasion has no purchase (Section 6.3) --- and invokes deeper semantic reading only behind that anchor. If any configuration escapes the trilemma, the data points here; building and pre-registering that rung against a fresh corpus is the natural next experiment.


\label{end-of-body}
\nocite{*}
\bibliographystyle{ACM-Reference-Format}
\bibliography{references}

\label{end-of-bib}
\appendix

\section{Generative AI Use Declaration}
Generative AI systems were used as instruments of the study and as drafting aids, in both cases under deterministic gating and human verification. As instruments: an LLM adversary proposed the out-of-vocabulary disguises of Section 6 (admitted only through a deterministic skeleton-invariance and novelty gate), and LLM judges (Sonnet, GPT-4o, Gemini) rated meaning-preservation in Sections 6 and 7 (trusted only after passing calibration controls). All such uses are logged in the closed run artifacts the paper is reproducible from. As a drafting aid: the authors used an LLM coding assistant to scaffold the figure renderer and verification gates and to draft prose, with every load-bearing number locked to a source artifact by an automated number-check and every claim verified by the authors. No citation, number, or result in this paper was generated without a verifiable source artifact.


\label{end-of-appendix}
\end{document}
